\subsection{Decidable constructive core}

\begin{theorem}[Jeandel--Rao eleven-tile witness as decidable constructive core]\label{thm:jeandel-rao-decidable-core}
  \lean{Obsmetrologia.AperiodicOrder.Tiling.Witness.ExistenceQuestion}
  \lean{Obsmetrologia.AperiodicOrder.Tiling.Witness.ExistenceQuestion.Assertion}
  \lean{Obsmetrologia.AperiodicOrder.Tiling.Witness.DecidableCore}
  \lean{Obsmetrologia.AperiodicOrder.Tiling.Witness.DecidableCore.ofCertificate}
  \lean{Obsmetrologia.AperiodicOrder.Tiling.Witness.DecidableCore.assertion}
  \lean{Obsmetrologia.AperiodicOrder.Tiling.Witness.jeandelRao_witness_is_decidable_core}
  \leanok
  \uses{thm:jeandel-rao-witness, thm:jeandel-rao-b3, thm:jeandel-rao-substitution-tower}
The finite Jeandel--Rao object with tile count \(11\), decorated cochain
\(\varphi\), and marker-ring core \(b_3\) is a decidable core of the
constructive aperiodic existence question.  Its certificate consists exactly of
the green finite nonzero class, the green marker-ring nonzero class, and the
green substitution-tower direct-limit persistence statement.

This is the precise green bridge behind the comparison with the domino problem:
an undecidable existence problem may only assert that some aperiodic object
exists, but this witness supplies the object together with a finite certificate
that the library can decide because the certificate is already proved.  The
node does not turn the domino problem into a decidable problem, and it does not
close the remaining identification with the literal Labb\'e recognizable
substitution matrix recorded at \cref{thm:jeandel-rao-persistence}.
\end{theorem}
